\clearpage\section{Centralbevægelse}\label{sec:central}
En central kraft har retning langs radius, $\vv F=F(r)\hat{\vv r}$, og $\vv L=\vv r\times\vv p$ er bevaret. Bevægelsen ligger derfor i et plan (for $\ell\ne0$). Adskil altid \emph{generelle centralpotentialer} fra de særlige formler for Keplerpotentialet $-k/r$.

\subsection{Fra to legemer til radial bevægelse}
For to masser $m_1,m_2$ defineres $M=m_1+m_2$, $\mu=m_1m_2/M$, $\vv R=(m_1\vv r_1+m_2\vv r_2)/M$ og $\vv r=\vv r_1-\vv r_2$. Positionerne kan rekonstrueres fra $\mathbf r_1=\mathbf R+(m_2/M)\mathbf r$ og $\mathbf r_2=\mathbf R-(m_1/M)\mathbf r$ for konventionen $\mathbf r=\mathbf r_1-\mathbf r_2$. Blandede led i den kinetiske energi ophæver hinanden.
For et internt potentiale $V(r)$ fås
\[T=\tfrac12M\dot{\vv R}^{2}+\tfrac12\mu\dot{\vv r}^{2},\qquad \mathcal L_{\rm rel}=\tfrac12\mu(\dot r^{2}+r^{2}\dot\phi^{2})-V(r).\]
Da $\phi$ er cyklisk, gælder
\[\boxed{\ell=\frac{\partial\mathcal L}{\partial\dot\phi}=\mu r^{2}\dot\phi=\mathrm{konstant}},\qquad\frac{\mathrm dA}{\mathrm dt}=\frac{\ell}{2\mu}.\]
For én partikel om et fast kraftcentrum erstattes $\mu$ af dens masse $m$.

\subsection{Det effektive potentiale: hvad bestemmer det?}\label{sec:veff}
\key{E=\frac12\mu\dot r^2+V_{\rm eff}(r),\qquad V_{\rm eff}(r)=V(r)+\frac{\ell^2}{2\mu r^2}.}
Den radiale bevægelse svarer til én dimension med et effektivt potentiale. Centrifugalleddet er ikke et nyt fysisk potentiale, men kommer fra den bevarede vinkelbevægelse.
\begin{center}\small\renewcommand{\arraystretch}{1.22}\begin{tabular}{p{47mm}p{91mm}}\toprule
\textbf{Spørgsmål} & \textbf{Matematisk test og fortolkning}\\\midrule
Hvor kan partiklen være? & $E\ge V_{\rm eff}(r)$, da $\mu\dot r^2/2\ge0$.\\
Hvor vender radialbevægelsen? & $E=V_{\rm eff}(r)$, dvs. $\dot r=0$; undersøg om punktet faktisk nås.\\
Hvor er en cirkulær bane? & $V_{\rm eff}'(r_0)=0$ med $r_0>0$.\\
Er cirkelbanen stabil? & $V_{\rm eff}''(r_0)>0$ giver radielt minimum; $<0$ giver ustabilt maksimum.\\
Er banen radielt begrænset? & Den tilladte radiale bevægelse er indesluttet i et endeligt interval; undersøg også randpunkter og singulariteter.\\
Er en begrænset bane lukket? & Kræver også $\Delta\phi_{\rm radial}/(2\pi)\in\mathbb Q$; radial begrænsning alene er ikke nok.\\\bottomrule
\end{tabular}\end{center}
\begin{center}\begin{tikzpicture}\begin{axis}[width=.77\linewidth,height=5cm,axis lines=left,xmin=.45,xmax=6.0,ymin=-.63,ymax=1.6,samples=180,domain=.50:6.0,xlabel={$r$},ylabel={$V_{\rm eff}(r)$},xtick=\empty,ytick=\empty,clip=false]
\addplot[accent,thick]{1/(2*x*x)-1/x};\addlegendentry{$-k/r+\ell^2/(2\mu r^2)$ (skaleret)}
\addplot[black!65,dashed,domain=.55:5.9]{-.18};\addlegendentry{$E$ (eksempel)}
\addplot[only marks,mark=*,mark size=1.5pt] coordinates {(.556,-.18) (5,-.18)};
\end{axis}\end{tikzpicture}\\[-1mm]{\small \itshape Et Kepler-eksempel: skæringer med energilinjen giver de radiale vendepunkter. Figuren er skematisk og dimensionsløs.}\end{center}

\subsection{Cirkulære baner og stabilitet}\label{sec:circular}
Radialligningen er
\[\mu\ddot r=\mu r\dot\phi^2-V'(r)=\frac{\ell^2}{\mu r^3}-V'(r)=-V_{\rm eff}'(r).\]
En cirkelbane har konstant $r=r_0$ og opfylder $V_{\rm eff}'(r_0)=0$. Sæt $r=r_0+\eta$, hvor $|\eta|\ll r_0$. Til første orden gælder
\[\mu\ddot\eta=-V_{\rm eff}''(r_0)\eta,\qquad\boxed{\omega_r^2=\frac{V_{\rm eff}''(r_0)}{\mu}}.\]
\begin{method}Givet $V(r)$: (1) tilføj $\ell^2/(2\mu r^2)$; (2) løs $V_{\rm eff}'(r_0)=0$; (3) indsæt $r_0$ i $V_{\rm eff}''(r_0)$ og undersøg fortegnet. Husk at kontrollere $r_0>0$ og parameterbetingelser.\end{method}
\textbf{Potenspotentiale:} Hvis $V=-a/r^{\alpha}$ med $a>0$, fås fra $V_{\rm eff}'(r_0)=0$
\[\frac{\ell^2}{\mu}=a\alpha r_0^{2-\alpha},\qquad V_{\rm eff}''(r_0)=\frac{\ell^2}{\mu r_0^4}(2-\alpha).\]
Altså er en eksisterende cirkelbane lineært stabil for $0<\alpha<2$ og ustabil for $\alpha>2$ (for $\alpha=2$ kræves særskilt analyse). Dette er \emph{ikke} en Kepler-specifik excentricitetsregel.

\subsection{Binet: fra $r(t)$ til $u(\phi)$}\label{sec:binet}
Sæt $u(\phi)=1/r$ og lad prim betegne afledning efter $\phi$, punkt efter $t$. Under bevaret $\ell\ne0$:
\begin{center}\small\renewcommand{\arraystretch}{1.22}\begin{tabular}{ll}\toprule Størrelse & Omskrivning\\\midrule
$r,\ r^2$ & $1/u,\ 1/u^2$\\
$\dot\phi$ & $\ell u^2/\mu$\\
$\dot r$ & $-(\ell/\mu)u'$\\
$\ddot r$ & $-(\ell^2u^2/\mu^2)u''$\\
$V(r),\ F(r)$ & $V(1/u),\ F(1/u)$\\\bottomrule\end{tabular}\end{center}
Fra kædereglen $\dot r=(\mathrm dr/\mathrm d\phi)\dot\phi$ får man $\dot r=-\ell u'/\mu$. Indsæt i $\mu(\ddot r-r\dot\phi^2)=F(r)$:
\key{u''+u=-\frac{\mu}{\ell^2u^2}F(1/u),\qquad r(\phi)=\frac{1}{u(\phi)}.}
Samtidig er energiudtrykket
\[\boxed{E=\frac{\ell^2}{2\mu}\big[(u')^2+u^2\big]+V(1/u)}.\]
\begin{method}Ud fra en radialkraft: find $F(1/u)$, opstil Binet-ODE'en, klassificér differentialligningen (harmonisk, eksponentiel osv.), løs $u(\phi)$, og brug energien eller begyndelsesdata til at bestemme integrationskonstanter.\end{method}

\subsection{Keplerbaner: energi, excentricitet og banetype}\label{sec:kepler}
\textbf{Kun for} $V(r)=-k/r$, $k>0$, gælder
\[u''+u=\frac{\mu k}{\ell^2},\qquad\boxed{r(\phi)=\frac{p}{1+e\cos(\phi-\phi_0)},\quad p=\frac{\ell^2}{\mu k}},\qquad\boxed{e=\sqrt{1+\frac{2E\ell^2}{\mu k^2}}}.\]
\begin{center}\small\renewcommand{\arraystretch}{1.18}\begin{tabular}{lll}\toprule Energi $E$ & Excentricitet $e$ & Banetype\\\midrule
$E<0$ & $0\le e<1$ & Ellipse (cirkel ved $e=0$)\\
$E=0$ & $e=1$ & Parabel\\
$E>0$ & $e>1$ & Hyperbel\\\bottomrule\end{tabular}\end{center}
\begin{center}\begin{tikzpicture}[scale=.82]
\draw[->,gray!70] (-3.35,0)--(3.4,0) node[right] {$x$};\draw[->,gray!70] (0,-2.1)--(0,2.25) node[above] {$y$};
\draw[accent,very thick] (-1.85,0) ellipse (1.4 and 1.02);
\fill[black] (-1.3,0) circle(1.6pt) node[above right,scale=.75]{Fokus};
\node[accent,scale=.85] at (-2.8,1.45) {Ellipse $e<1$};
\begin{scope}[xshift=2.05cm,yshift=-.05cm]
\draw[red!65!black,very thick,domain=-1.18:1.18,samples=90] plot ({.55+0.48*\x*\x},{1.05*\x});
\node[red!65!black,scale=.85] at (1.0,1.65) {Parabel $e=1$};
\end{scope}
\end{tikzpicture}\\[-1mm]{\small\itshape Skematisk illustration af keglesnit; hyperblen svarer til $e>1$.}\end{center}
\textbf{Ellipsegeometri:} Kraftcentret ligger i et brændpunkt, ikke i centrum. For en ellipse gælder $r_{\min}=a(1-e)$, $r_{\max}=a(1+e)$ og $b=a\sqrt{1-e^2}$.
\begin{center}\begin{tikzpicture}[scale=.82]
\draw[->,gray!60](-3.6,0)--(3.65,0) node[right]{$x$};
\draw[accent,thick](0,0) ellipse (3 and 1.65);
\fill[gray!65](0,0) circle (1.5pt) node[above right,scale=.77]{Centrum};
\fill[black](-2.5,0) circle (1.7pt) node[below,scale=.77]{Fokus};
\draw[accent,->,thick](-2.5,0)--(-3,0) node[midway,above,scale=.7]{$r_{\min}$};
\draw[accent,->,thick](-2.5,-.13)--(3,-.13) node[midway,below,scale=.7]{$r_{\max}$};
\node[scale=.8] at (-3.0,.33){Periapsis};\node[scale=.8] at (2.5,.33){Apoapsis};
\end{tikzpicture}\\{\small\itshape Keplerellipsens fokus ligger forskudt fra ellipsens centrum.}\end{center}
\textbf{Ellipseformler:} $E=-k/(2a)$, $r_{\min}=a(1-e)$, $r_{\max}=a(1+e)$, og Keplers tredje lov er $T^2=4\pi^2\mu a^3/k$. Arealloven $\dot A=\ell/(2\mu)$ gælder derimod for enhver central kraft.

\subsection{Udledning og kontrol af excentriciteten}
Med $V=-ku$ og $\dot r=-\ell u'/\mu$ er
\[E=\frac{\ell^2}{2\mu}[(u')^2+u^2]-ku.\]
Sæt $u=(\mu k/\ell^2)(1+e\cos\phi)$ (periapsis ved $\phi=0$):
\[E=\frac{\mu k^2}{2\ell^2}\big[e^2\sin^2\phi+(1+e\cos\phi)^2-2(1+e\cos\phi)\big]=\frac{\mu k^2}{2\ell^2}(e^2-1).\]
Det giver excentricitetsformlen ovenfor. Ellipsens areal $\pi ab$ og $b^2=ap$ sammen med arealhastigheden giver $T^2=4\pi^2\mu a^3/k$.

\subsection{Radialperiode, lukning og præcession}
Hvis der findes to radiale vendepunkter, er vinkelændringen mellem dem
\[\Delta\phi_{1/2}=\int_{r_{\min}}^{r_{\max}}\frac{\ell\,\mathrm dr}{r^2\sqrt{2\mu[E-V_{\rm eff}(r)]}},\qquad\Delta\phi_{\rm radial}=2\Delta\phi_{1/2}.\]
En bane er lukket, når $\Delta\phi_{\rm radial}/(2\pi)$ er rationalt. Ved Keplerpotentialet er den $1$, mens andre centralpotentialer kan give præcession. Forskydning fra én periapsis til den næste er $\Delta\phi_{\rm radial}-2\pi$ ved sammenligning med én fuld omdrejning.

\subsection{Modificeret Keplerpotentiale og præcession}
For $F(r)=-k/r^2+\gamma/r^3$ med $k>0$ svarer potentialet til $V(r)=-k/r+\gamma/(2r^2)$. Binet giver
\[u''+\left(1+\frac{\mu\gamma}{\ell^2}\right)u=\frac{\mu k}{\ell^2}.\]
Sæt $\nu^2=1+\mu\gamma/\ell^2>0$. En oscillerende løsning er
\[\boxed{r(\phi)=\frac{\ell^2\nu^2/(\mu k)}{1+e\cos[\nu(\phi-\phi_0)]}},\qquad
\Delta\phi_{\rm radial}=\frac{2\pi}{\nu}.\]
For $\nu\ne1$ er den næste periapsis generelt forskudt i forhold til én fuld omdrejning; banen er kun lukket for rationelle forhold mellem radial- og vinkelperioderne. Ved en næsten cirkulær bane er $\Delta\phi_{\rm radial}\simeq2\pi\omega_\phi/\omega_r$, hvor $\omega_\phi=\ell/(\mu r_0^2)$.

\subsection{Scattering og virialsætning}
Ved asymptotisk hastighed $v_\infty$ og anslagsparameter $b$ er $E=\tfrac12\mu v_\infty^2$, $\ell=\mu v_\infty b$. For en enkeltvendepunktsbane kan afbøjningsvinklen findes fra
\[\phi_0=\int_{r_{\min}}^\infty\frac{\ell\,\mathrm dr}{r^2\sqrt{2\mu[E-V_{\rm eff}(r)]}},\qquad\Theta=\pi-2\phi_0,\]
med valgt fortegnskonvention. Det differentielle tværsnit er $\mathrm d\sigma/\mathrm d\Omega=(b/\sin\Theta)|\mathrm db/\mathrm d\Theta|$ for én gren; summer over grene, hvis flere $b$ giver samme $\Theta$.
For en bundet Keplerbane med $V=-k/r$ følger særligt $\langle T\rangle=-E$ og $\langle V\rangle=2E$, hvor $E<0$. Virialen $G=\mathbf r\cdot\mathbf p$ er i sig selv generelt ikke bevaret: $G>0$ når afstanden vokser og $G<0$ når den aftager.

For en bundet bevægelse med veldefineret langtidsgennemsnit sættes $G=\vv r\cdot\vv p$, så
\[\frac{\mathrm dG}{\mathrm dt}=2T-\vv r\cdot\nabla V,\qquad2\langle T\rangle=\langle\vv r\cdot\nabla V\rangle.\]
Ved homogent $V\propto r^n$ giver det $2\langle T\rangle=n\langle V\rangle$.

\subsection{Laplace--Runge--Lenz-vektoren}
\begin{center}\begin{tikzpicture}[scale=.72]
\draw[accent,thick](0,0) ellipse (3 and 1.6);\fill[black](-2.4,0) circle (2pt) node[below,scale=.8]{Kraftcentrum};
\draw[->,red!75!black,very thick](-2.4,0)--(-3.1,0) node[midway,above,scale=.85]{$\mathbf A$};
\node[scale=.8] at (-2.9,.42){Periapsis};
\end{tikzpicture}\\{\small\itshape Den bevarede Runge--Lenz-vektor peger mod periapsis og fastlægger retningen.}\end{center}
For $V=-k/r$ er den bevarede vektor $\vv A=\vv p\times\vv L-\mu k\hat{\vv r}$ med $|\vv A|=\mu k e$. Da $\dot{\vv p}=-k\vv r/r^3$ og $\dot{\vv L}=0$, fås ved tredobbelt krydsprodukt
\[\dot{\vv p}\times\vv L=\mu k\left(\frac{\dot{\vv r}}r-\frac{\vv r(\vv r\cdot\dot{\vv r})}{r^3}\right)=\mu k\frac{\mathrm d\hat{\vv r}}{\mathrm dt},\]
så $\dot{\vv A}=0$. Desuden er $\vv A\cdot\vv r=\ell^2-\mu kr$, som giver banens polarform med periapsis i retning af $\vv A$.

